\section{Khó khăn và giải pháp}

Trong quá trình thiết kế và triển khai hệ thống, nhóm gặp phải nhiều thách thức kỹ thuật liên quan đến lập trình socket,
đồng bộ hóa luồng và truyền dữ liệu. Phần này trình bày các vấn đề thực tế được ghi nhận trong mã nguồn,
phân biệt rõ giữa những vấn đề đã được giải quyết và những vấn đề còn tồn tại.

% =========================================================
\subsection{Các vấn đề đã được giải quyết}
% =========================================================

\subsubsection{Phân mảnh gói tin TCP}

TCP là giao thức hướng luồng (stream-oriented): không có ranh giới gói tin tự nhiên, dữ liệu có thể đến theo từng mảnh nhỏ hoặc bị ghép lại tùy ý.
Nếu không có cơ chế phân tách gói tin rõ ràng, bên nhận sẽ không thể xác định điểm bắt đầu và kết thúc của mỗi thông điệp.

\textbf{Giải pháp:} \texttt{PacketCodec} đọc độ dài 4 byte big-endian rồi đọc đủ
JSON UTF-8. Luồng file đọc đúng số byte khai báo trước khi quay lại đọc frame
\texttt{FILE\_END}; việc TCP chia nhỏ hoặc gộp dữ liệu không làm mất ranh giới.

\subsubsection{Đồng bộ hóa ghi lên socket dùng chung}

Nhiều virtual thread có thể cùng lúc muốn ghi lên cùng một \texttt{OutputStream} của một socket.
Nếu không đồng bộ hóa, các byte từ hai gói tin khác nhau sẽ bị xen kẽ, gây lỗi phân tích giao thức ở phía nhận.

\textbf{Giải pháp:} Mọi nơi ghi lên một socket đều đi qua cùng \texttt{PacketChannel}.
Khóa ghi bảo vệ thao tác cấp sequence và ghi trọn frame, bao gồm phản hồi,
broadcast và heartbeat. Đóng socket không chờ khóa ghi để giải phóng writer đang bị kẹt.

\subsubsection{Phân luồng kết nối đến trên cùng một cổng}

Server lắng nghe trên một cổng TCP duy nhất (8080). Tuy nhiên, cổng này phải phục vụ ba loại kết nối khác nhau:
phiên điều khiển thông thường, luồng upload file broadcast, và luồng download file broadcast.
Tương tự, \texttt{PeerListener} phía client cũng phải phân biệt kết nối chat và kết nối dữ liệu file.

\textbf{Giải pháp:} Cả hai đều sử dụng chiến lược \textit{đọc gói tin đầu tiên để phân loại}.
Tại \texttt{ClientHandler.run()}: nếu gói đầu tiên là \texttt{HELLO} → khởi chạy phiên điều khiển;
nếu là \texttt{FILE\_BROADCAST\_UPLOAD\_START} → chuyển sang trình xử lý upload;
nếu là \texttt{FILE\_BROADCAST\_DOWNLOAD\_START} → chuyển sang trình xử lý download.
Tại \texttt{PeerListener}: nếu là \texttt{PEER\_HELLO} → khởi tạo kết nối chat;
nếu là \texttt{FILE\_OFFER} → chuyển cho \texttt{FileTransferController} xử lý luồng file.

\subsubsection{Điều kiện tranh chấp trong thiết lập kết nối P2P}

Khi hai client cùng thực hiện \texttt{/lookup} nhau gần như đồng thời, có thể xảy ra trường hợp cả hai đều gửi yêu cầu kết nối,
dẫn đến hai kết nối P2P song song giữa cùng một cặp client.

\textbf{Giải pháp:} Mỗi kết nối chat mang một UUID. Khi hai handshake giao nhau,
cả hai phía giữ kết nối có UUID nhỏ hơn; map peer được cập nhật dưới khóa của
\texttt{ClientSession}. Callback đóng của kết nối cũ chỉ được xóa chính kết nối đó,
không được xóa socket thay thế. Báo cáo metadata được xếp hàng riêng theo thứ tự.

\subsubsection{Tính toàn vẹn dữ liệu file}

File truyền qua mạng có thể bị hỏng do lỗi đường truyền hoặc gián đoạn kết nối.
Ngoài ra, file tạm thời (đang ghi) không được để lộ ra hệ thống tệp cho đến khi quá trình nhận hoàn tất.

\textbf{Giải pháp:} \texttt{FileSender} tính toán SHA-256 trực tiếp trong quá trình gửi chunk và đính kèm vào gói \texttt{FILE\_END}.
\texttt{FileReceiver} tính lại SHA-256 độc lập trong quá trình nhận và so sánh với giá trị nhận được.
Nếu khớp, file tạm được di chuyển đến vị trí đích bằng \texttt{Files.move(..., StandardCopyOption.ATOMIC\_MOVE)};
nếu không khớp, file tạm bị xóa và trả về \texttt{FILE\_RESULT(CHECKSUM\_MISMATCH)}.

% =========================================================
\subsection{Vấn đề còn tồn tại}
% =========================================================

\subsubsection{Không có cơ chế bền vững (Persistence)}

Toàn bộ trạng thái hệ thống — bao gồm phiên đăng nhập, danh sách cộng đồng, thành viên, và metadata file broadcast —
chỉ tồn tại trong bộ nhớ (\texttt{ConcurrentHashMap} trong RAM).
Khi server khởi động lại, mọi dữ liệu đều bị mất.
\texttt{AuthService} hiện tại sử dụng tài khoản hardcode, không có cơ sở dữ liệu hay file cấu hình.
Đây là giới hạn thiết kế quan trọng cần giải quyết trước khi triển khai trong môi trường thực tế.

\subsubsection{Giới hạn phục hồi kết nối}

Heartbeat hai chiều đã phát hiện thiếu ACK trên control socket. Người dùng có thể
kết nối lại bằng mật khẩu và tải trạng thái mới; P2P còn sống được giữ nguyên.
App chưa tự nối lại một peer đã chết và chưa resume file theo offset; file bị ngắt
phải gửi lại. Thử nghiệm localhost chưa thay thế được kiểm thử nhiều máy qua NAT/firewall.
